\label{inv:A3.2}\label{wb:lem:bcpr-eta-extension}
\StatementRecord{inv:A3.2}{proved}{}{the hypotheses, notation and definitions stated here; \cref{inv:A2.4}; \cref{inv:eq:A-E-line}; \cref{inv:A1.3}}

The Leray associated-graded spaces for the four middle groups are
\[
\begin{array}{c|cc}
&E_2^{1,q-1}&E_2^{0,q}\\\hline
H^{1,1}&H^1(\OO(-2))=\C\alpha&\C\eta_{\rm rel}\\
H^{1,2}&0&\C\upsilon\\
H^{2,1}&\C\upsilon^\vee&0\\
H^{2,2}&\C\eta_{\rm rel}^\vee&\C\alpha^\vee.
\end{array}
\]
Here $\alpha$ is the Cech class of $dt/t$ on the two standard affine
base charts, $\eta_{\rm rel}=-q_E$ is the invariant relative class
from \ref{inv:A2.4}, and $\upsilon$ is the image of the trace-dual of
$1\in f_*\OO(H)$ under $R^2f_*N\to R^2f_*\mathcal E$ and the
isomorphism $R^2f_*\Omega_X^1=R^2f_*\mathcal E$.
The superscript $\vee$ denotes the bases dual under the relative
trace and the base Serre pairing.  There exists a Dolbeault class
$Q\in H^1(X,\Omega_X^1)$ whose restriction to every regular fibre
is $\eta=u\wedge w+6\gamma\wedge\delta$.
\end{proposition}


\begin{proof}
\begin{enumerate}[label=\arabic*.,leftmargin=*]
\item The isomorphism $f_*\Omega_X^1=\Omega_B^1$ in the residue sequence
is the actual pullback map, so $\alpha$ gives its stated Leray
class; it can be represented by a nonzero normalization of
$f^*\omega_{\rm FS}$ in Dolbeault cohomology.  The section
$\eta_{\rm rel}$ belongs to the canonical kernel identification
$R^1f_*\Omega_X^1=\ker\kappa_1$ and spans its global sections.
Leray is surjective onto this space, so any preimage defines $Q$
with the asserted fibre restriction.  Its ambiguity is exactly
$\C\alpha$; the statement does not claim a canonical global form.
The injection $R^2f_*N=\OO\to R^2f_*\mathcal E$ from
\eqref{inv:eq:A-E-line} spans its global sections, since the quotient
is $\OO(-1)$.  This defines $\upsilon$ with no choice of a negative
degree summand.  Duality used in \ref{inv:A1.3} gives the remaining entries
and their specified dual pairings.  Choosing a lift of
$\eta_{\rm rel}$ and of $\alpha^\vee$ gives bases of the whole
middle cohomology groups.  The associated-graded descriptions are
canonical with the named choices; the splitting of Leray vector
spaces is not asserted canonical.
\end{enumerate}
\end{proof}